\chapter{ValDX}
\label{ch:valdx}

%  The structure follows the manuscript’s results argument (projects/1-ValDX/manuscript/report/7_results_omc.tex:4).

%    Section                                                                Purpose and material
%   ━━━━━━━━━━━━━━━━━━━━━━━━━━━━━━━━━━━━━━━━━━━━━━━━━━━━━━━━━━━━━━━━━━━━━  ━━━━━━━━━━━━━━━━━━━━━━━━━━━━━━━━━━━━━━━━━━━━━━━━━━━━━━━━━━━━━━━━━━━━━━━━━━━━━━━━━━━━━━━━━━━━━━━━━━━━━━━━━━━━━━━━━━━━━━
%    1. Introduction and objectives                                         Briefly establish the unresolved problem inherited from the background chapter. Similar uptake agreement can support
%                                                                           different structural populations. Introduce held-out prediction, modification from the prior and replicate
%                                                                           consistency as complementary assessments.
%   ─────────────────────────────────────────────────────────────────────  ──────────────────────────────────────────────────────────────────────────────────────────────────────────────────────
%    2. Validation framework and methods                                    Give the workflow and definitions one home. Explain overlap-aware partitions, replicate fits, training-only
%                                                                           selection of fitting strength, consensus construction, Apparent Work and the four ValDX Work metrics. Include
%                                                                           compact tables of datasets, candidate ensembles and fitting protocols.
%   ─────────────────────────────────────────────────────────────────────  ──────────────────────────────────────────────────────────────────────────────────────────────────────────────────────
%    3. Uptake agreement and population recovery                            Retain TeaA as the opening result. Known populations establish what training error misses and what Apparent Work
%                                                                           adds. Preserve the distinction between increasing Work along one fitting trajectory and comparing Work across
%                                                                           candidates at comparable agreement.
%   ─────────────────────────────────────────────────────────────────────  ──────────────────────────────────────────────────────────────────────────────────────────────────────────────────────
%    4. Comparing structural hypotheses                                    Present BPTI reweighting, then HOIP parameter fitting as two subsections. BPTI tests ensemble comparison using NMR-
%                                                                           derived residue-wise uptake; HOIP extends assessment to individual structures and exposes the limits of ensemble-
%                                                                           average parameter fitting.
%   ─────────────────────────────────────────────────────────────────────  ──────────────────────────────────────────────────────────────────────────────────────────────────────────────────────
%    5. Population fitting and forward-model calibration                    Place the BPTI protocol comparison directly beside the kinetic-misspecification benchmark. Their different preferred
%                                                                           orders make the governing conclusion explicit: fitting order depends on the mismatch being addressed. Distinguish
%                                                                           final-stage Work from total modification across stages.
%   ─────────────────────────────────────────────────────────────────────  ──────────────────────────────────────────────────────────────────────────────────────────────────────────────────────
%    6. Clustering and diagnostic perturbations                            Promote the clustering sweep and BRD4 artefact experiments from the supplement. Clustering tests whether conclusions
%                                                                           survive reduced representation; perturbations test sensitivity to specific structural and forward-model inputs. Keep
%                                                                           their main figures and shorten the exhaustive comparisons.
%   ─────────────────────────────────────────────────────────────────────  ──────────────────────────────────────────────────────────────────────────────────────────────────────────────────────
%    7. Discussion and chapter conclusions                                  Draw out practical decisions from the combined evidence. End on the remaining limits—fixed candidate coverage,
%                                                                           forward-model suitability and uncertainty whose cause remains unresolved—and use these to motivate jaxENT.

%   Keep general HDX kinetics, BV theory and MaxEnt foundations in the background chapter. Move detailed generation settings, simulation protocols, digestion parameters, simulator checks and
%   extended tables into the existing ValDX protocols appendix (thesis/appendices/a-valdx-protocols.tex:1). Reuse the manuscript’s mature results prose and figures, editing transitions and
%   repeated definitions.

%   The BICePs comparison already exists in the revised methods (projects/1-ValDX/manuscript/report/6_methods_omc.tex:464). Keep a short subsection distinguishing Bayesian evidence from ValDX’s
%   diagnostic distances. A direct benchmark would be additional work.

% in this chapter we introduct data splitting, methods for data splitting, and interpretable statistical metrics used for comparative analysis. The insight from these metrics are that considering differences in ensemble-models that reflect the phenoma observed during gradient-based optimization of the ensemble-models. Ie that weights and parameters should be considered jointly. We compare this to BiCePs score, which is a complete bayesian formulation.
% in valdx we employ our approach across a range of experiments: ensemble-comparison, single-structure comparison, and joint inference of populations and forward-model parameters. We also show that clustering and decoy conformations can be used to probe ensemble qualites and representativeness.

% We need to define the split types and how they are constructed.

% We need to define the work metrics as well as BiCePs score and show how they are related/different.

\section{Introduction and objectives}
\label{sec:valdx-intro}

\section{Validation framework and methods}
\label{sec:valdx-methods}

\subsection{Experimental data and candidate ensembles}
\label{sec:valdx-inputs}
% Dataset provenance, preprocessing overview, and ensemble hypotheses.
% Distinguish BPTI NMR-derived uptake from measured peptide uptake.

\subsection{Data splitting and replicate fits}
\label{sec:valdx-splits}
% Baseline, Non-Redundant and Spatial partitions.
% Overlap removal, training fractions, and replicate design.

\subsection{Fitting protocols, parameter selection and consensus}
\label{sec:valdx-fitting}
% RW-only, BV-only, staged and alternating fitting.
% Training-only selection of fitting strength and consensus construction.

\subsection{Assessment metrics}
\label{sec:valdx-work}
% Uptake error, population recovery, Apparent Work and ValDX Work metrics.
% Replicate dispersion, agreement between partitions, and Work accounting.
% Brief comparison with Bayesian evidence scoring.

\subsection{Benchmark designs and statistical analysis}
\label{sec:valdx-benchmarks}

\subsubsection{Synthetic population and kinetic benchmarks}
\label{sec:valdx-synthetic}
% TeaA target populations and candidate composition.
% Forward-model mismatch and controlled changes in exchange kinetics.

\subsubsection{Clustering and structural perturbations}
\label{sec:valdx-decoys}
% Clustering representation and size sweep.
% Gaussian noise, coordinate mixing and proton shuffling.

\subsubsection{Statistical analysis}
\label{sec:valdx-statistics}
% Between-candidate comparisons, changes against zero,
% regression, split sensitivity and descriptive analyses.

\section{Uptake agreement and population recovery}
\label{sec:valdx-population-recovery}

\section{Comparing structural hypotheses}
\label{sec:valdx-structural-hypotheses}

\subsection{BPTI: comparing ensembles through reweighting}
\label{sec:valdx-bpti-reweighting}

\subsection{HOIP: comparing ensembles and individual structures through BV parameter fitting}
\label{sec:valdx-hoip-parameter-fitting}

\section{Population fitting and forward-model calibration}
\label{sec:valdx-calibration}

\subsection{BPTI: comparison of fitting protocols}
\label{sec:valdx-protocol-comparison}

\subsection{TeaA: response to kinetic misspecification}
\label{sec:valdx-kinetic-misspecification}

\subsection{Fitting order and the source of mismatch}
\label{sec:valdx-fitting-order}

\section{Clustering and diagnostic perturbations}
\label{sec:valdx-perturbations}

\subsection{Sensitivity to ensemble clustering}
\label{sec:valdx-clustering-results}

\subsection{BRD4 structural and forward-model perturbations}
\label{sec:valdx-brd4-perturbations}

\section{Discussion and chapter conclusions}
\label{sec:valdx-discussion}
